\section{\MakeUppercase{Performance Evaluation and Analysis}}\label{chap:three}
This chapter reports an experimental assessment of the proposed audio classification frameworks across varied acoustic settings, spanning industrial, urban, and natural soundscapes. Several public datasets are used to evaluate model robustness and generalization under heterogeneous, noisy conditions. Although anomaly event detection remains the core objective of this dissertation, benchmark datasets such as \gls{esc50} and \gls{fsc22}, which are typically framed as multi-class classification problems, are instead used as proxy testbeds to investigate anomaly-related characteristics, including resilience to background noise, detection of rare or short-duration events, and consistency across diverse sound classes.

\dots

% You need this part to separate between English and Lithuanian Part

\addtocontents{lof}{\protect\vspace*{2em}}
\addtocontents{lof}{\textbf{\MakeUppercase{Paveikslų sąrašas}}\par}
\addtocontents{lof}{\protect\vspace*{1em}}

\addtocontents{lot}{\protect\vspace*{2em}}
\addtocontents{lot}{\textbf{\MakeUppercase{Lentelių sąrašas}}\par}
\addtocontents{lot}{\protect\vspace*{1em}}
